\documentclass[11pt]{article}
\usepackage[T1]{fontenc}
\usepackage{textcomp}
\usepackage[margin=0.75in]{geometry}
\usepackage{amsmath,amssymb,amsthm}
\usepackage{array,booktabs}
\usepackage{hyperref}
\setlength{\parindent}{0pt}
\newtheorem{theorem}{Theorem}
\theoremstyle{definition}
\newtheorem{definition}{Definition}
\title{Flow Proof Artifact: prop-11-perpendicular-at-point}
\author{Generated by \texttt{flow doc proof}}
\date{Flow Proof Book}
\begin{document}
\maketitle
To draw a straight line at right angles to a given straight line from a given point on it.\\[0.5em]
\textbf{Source.} euclid — Elements, Book I, Proposition 11\\[0.5em]
\bigskip
\noindent{\large \textbf{Derived fact 1.} \textit{To draw a straight line at right angles to a given straight line from a given point on it} \hfill the Euclidean plane \cdot Euclid Book I \cdot Proposition 11: erect a perpendicular at a point on a line}\par
\smallskip\noindent\textit{Coordinate: the Euclidean plane \cdot Euclid Book I \cdot Proposition 11: erect a perpendicular at a point on a line}\par
\smallskip\noindent\textit{Source: euclid — Elements, Book I, Proposition 11}\par
\smallskip\noindent\textit{Built on: proposition 1: equilateral triangle on a segment, for Book I of the Elements in on the Euclidean plane, proposition 4: side-angle-side congruence, for Book I of the Elements in on the Euclidean plane, proposition 8: side-side-side angle equality, for Book I of the Elements in on the Euclidean plane}\par
\medskip
\noindent\textbf{Goal.} To draw a straight line at right angles to a given straight line from a given point on it\par
\noindent$\displaystyle \angle APB = 90^\circ$\par
\medskip
\noindent
\begin{tabular}{@{} >{\raggedright\arraybackslash}p{0.06\textwidth} >{\raggedright\arraybackslash}p{0.47\textwidth} >{\raggedleft\arraybackslash}p{0.41\textwidth} @{}}
\textbf{\#} & \textbf{Proof} & \textbf{Mathematics} \\
\midrule
\textbf{1.} & We prove that proposition 11: erect a perpendicular at a point on a line for Book I of the Elements in the Euclidean plane. & \\[0.45em]
\textbf{2.} & Let C = point\_on\_AB\_with\_AC\_equal\_to\_equilateral\_side. & \\[0.45em]
\textbf{3.} & Let equilateral\_ACD = triangle\_on\_AC\_by\_Proposition\_1. & \\[0.45em]
\textbf{4.} & Let equilateral\_BCE = triangle\_on\_BC\_by\_Proposition\_1. & \\[0.45em]
\textbf{5.} & Let P = intersection\_of\_segments\_from\_D\_and\_E\_above\_AB. & \\[0.45em]
\textbf{6.} & We invoke the axiom governing Book I of the Elements in the Euclidean plane: proposition 1: equilateral triangle on a segment, for Book I of the Elements in on the Euclidean plane. & \\[0.45em]
\textbf{7.} & We invoke the axiom governing Book I of the Elements in the Euclidean plane: proposition 4: side-angle-side congruence, for Book I of the Elements in on the Euclidean plane. & \\[0.45em]
\textbf{8.} & We invoke the derived fact governing Book I of the Elements in the Euclidean plane: proposition 8: side-side-side angle equality, for Book I of the Elements in on the Euclidean plane. & \\[0.45em]
\textbf{9.} & From step 2, step 3, step 4, step 5, step 6, step 7, and step 8, we can deduce that angle APC equals angle BPC. & $\displaystyle \angle APC = \angle BPC$ \\[0.45em]
\textbf{10.} & We can deduce that angle APB equals one right angle. Hence proven. & $\displaystyle \angle APB = 90^\circ$ \\[0.45em]
\bottomrule
\end{tabular}
\label{thm:Geometry-Euclid Book I-proposition_11_erect_a_perpendicular_at_a_point_on_a_line}
\par\medskip\hrule
\end{document}
